% ============================================================
%  part4/ch20-existence-uniqueness.tex
%  Chapter 20: Existence and Uniqueness
% ============================================================

\chapter{Existence and Uniqueness}
\label{ch:existence-uniqueness}

\section{The RSVP free-energy functional}

The basic RSVP state is
$X(t) = (\Phi, \mathbf{v}, S, \ell)$ with $L = \ell^2 \geq 0$.

\begin{definition}[RSVP free-energy functional]
  \label{def:rsvp-functional}
  \[
    \mathcal{F}[\Phi,\mathbf{v},S,\ell]
    = \int_\Omega \Bigl[
      \tfrac{\alpha}{2}|\nabla\Phi|^2 + V(\Phi)
      + \tfrac{\beta}{2}|\nabla\ell|^2 + U(\ell)
      + \lambda W(\Phi)\ell^2
      + \tfrac{\rho}{2}|\mathbf{v}|^2
      + \tfrac{\eta}{2}|\nabla\times\mathbf{v}|^2
      + \tfrac{\chi}{2}|\nabla S|^2 - TS
    \Bigr]dx.
  \]
\end{definition}

The coupling term $\lambda W(\Phi)\ell^2$ is bounded below
whenever $\lambda \geq 0$ and $W(\Phi) \geq 0$.
Natural choices: $W(\Phi) = \Phi^2$ or $W(\Phi) = |\nabla\Phi|^2$.

\section{Functional derivatives}

\begin{proposition}[Variational derivatives of $\mathcal{F}$]
  \label{prop:rsvp-variations}
  \statusproven{}
  Under periodic or no-flux boundary conditions:
  \begin{align}
    \frac{\delta\mathcal{F}}{\delta\Phi}
    &= -\alpha\Delta\Phi + V'(\Phi) + \lambda W'(\Phi)\ell^2,
    \label{eq:dFdPhi}\\
    \frac{\delta\mathcal{F}}{\delta\ell}
    &= -\beta\Delta\ell + U'(\ell) + 2\lambda W(\Phi)\ell,
    \label{eq:dFdell}\\
    \frac{\delta\mathcal{F}}{\delta S}
    &= -\chi\Delta S - T,
    \label{eq:dFdS}\\
    \frac{\delta\mathcal{F}}{\delta\mathbf{v}}
    &= \rho\mathbf{v} + \eta\nabla\times(\nabla\times\mathbf{v}).
    \label{eq:dFdv}
  \end{align}
\end{proposition}

\begin{proof}
  For $\Phi$: vary $\Phi \mapsto \Phi + \varepsilon\psi$ and
  differentiate at $\varepsilon = 0$.
  The gradient term gives $-\alpha\Delta\Phi$ after integration
  by parts.
  The coupling term gives $\lambda W'(\Phi)\ell^2$.
  Equations~\eqref{eq:dFdell}--\eqref{eq:dFdv} follow
  analogously.
\end{proof}

\section{RSVP gradient flow}

\begin{definition}[Pure dissipative RSVP system]
  \label{def:rsvp-gf}
  \begin{align*}
    \partial_t\Phi &= M_\Phi[\alpha\Delta\Phi - V'(\Phi) - \lambda W'(\Phi)\ell^2],\\
    \partial_t\ell &= M_\ell[\beta\Delta\ell - U'(\ell) - 2\lambda W(\Phi)\ell],\\
    \partial_t S   &= M_S(\chi\Delta S + T),\\
    \partial_t\mathbf{v} &= -M_v[\rho\mathbf{v} + \eta\nabla\times(\nabla\times\mathbf{v})].
  \end{align*}
\end{definition}

\begin{theorem}[RSVP free-energy dissipation]
  \label{thm:rsvp-dissipation}
  \statusproven{}
  For the pure gradient-flow RSVP system with
  $M_\Phi, M_\ell, M_S, M_v \geq 0$,
  \[
    \frac{d\mathcal{F}}{dt}
    = -M_\Phi\!\int\!\Bigl|\frac{\delta\mathcal{F}}{\delta\Phi}\Bigr|^2 dx
    - M_\ell\!\int\!\Bigl|\frac{\delta\mathcal{F}}{\delta\ell}\Bigr|^2 dx
    - M_S\!\int\!\Bigl|\frac{\delta\mathcal{F}}{\delta S}\Bigr|^2 dx
    - M_v\!\int\!\Bigl|\frac{\delta\mathcal{F}}{\delta\mathbf{v}}\Bigr|^2 dx
    \leq 0.
  \]
\end{theorem}

\begin{proof}
  By the chain rule for functionals,
  $\frac{d\mathcal{F}}{dt}
  = \sum_q \int \frac{\delta\mathcal{F}}{\delta q}\,\partial_t q\,dx$.
  Substituting $\partial_t q = -M_q \frac{\delta\mathcal{F}}{\delta q}$
  for each field $q \in \{\Phi, \ell, S, \mathbf{v}\}$ gives
  $-M_q \int |\frac{\delta\mathcal{F}}{\delta q}|^2 dx \leq 0$.
\end{proof}

\section{Conservative lamphron variant}

\begin{theorem}[Lamphron mass conservation]
  \label{thm:lamphron-mass}
  \statusproven{}
  For the Cahn--Hilliard-type lamphron evolution
  $\partial_t L = \nabla\cdot(M_L\nabla\mu_L)$
  with $\mu_L = \delta\mathcal{F}/\delta L$ and no-flux boundaries,
  \[
    \frac{d}{dt}\int_\Omega L\,dx = 0,
    \qquad
    \frac{d\mathcal{F}}{dt}
    = -\int_\Omega M_L|\nabla\mu_L|^2\,dx \leq 0.
  \]
\end{theorem}

\begin{proof}
  Mass conservation: integrate the evolution equation and apply
  the divergence theorem.
  Energy dissipation: $\frac{d\mathcal{F}}{dt}
  = \int\mu_L\partial_tL\,dx = M_L\int\mu_L\Delta\mu_L\,dx
  = -M_L\int|\nabla\mu_L|^2\,dx \leq 0$.
\end{proof}

\section{Reduced existence results and the global conjecture}

\begin{StatusBox}{Proven Framework}
  The following subsystems have classical existence theory:
  Allen--Cahn $\partial_t\Phi = M_\Phi[\alpha\Delta\Phi - V'(\Phi)]$;
  Cahn--Hilliard $\partial_t L = \nabla\cdot(M_L\nabla\mu_L)$.
\end{StatusBox}

\begin{conjecture}[Global weak RSVP existence]
  \label{conj:rsvp-existence}
  \statusconjectured{}
  Let $(\Phi_0, \ell_0, S_0, \mathbf{v}_0) \in H^k(\Omega)$
  with $k$ sufficiently large, finite free energy, and
  compatible boundary conditions.
  Then the coupled RSVP system admits a global weak solution
  satisfying $\frac{d\mathcal{F}}{dt} \leq 0$ in the
  distributional sense.
\end{conjecture}

\begin{OpenQuestion}[RSVP well-posedness]
  Determine the minimal Sobolev index $k$ for which the
  coupled RSVP system is well-posed, and prove or disprove
  Conjecture~\ref{conj:rsvp-existence}.
  The main difficulty is controlling the nonlinear coupling
  term $\lambda W(\Phi)\ell^2$ under advection.
\end{OpenQuestion}
